\begin{luacode*}
    dofile('./protocol.lua')

    function print_spectral_rows(data_key)
        local data = protocol[data_key]
        for i = 1, #data.Divisions do
            -- Считаем дополнительные параметры (пример)
            if data_key == "SpectralCdS" then
                k = 1e3
            else
                k = 1e6
            end
            local Rc_Mohm = data.R_c[i] / k
            
            tex.print(string.format(
                "$%d$ & $%.3f$ & $%.3f$ & $%.3f$ & & & & \\\\", 
                data.Divisions[i], data.lambda[i] * 1e6, data.E_lambda[i], Rc_Mohm
            ))
        end
    end

    function print_intensity_rows()
        local data = protocol.Intensity
        for i = 1, #data.R_c do
            local Rc_Mohm = data.R_c[i] / 1e6
            tex.print(string.format(
                "$%.2f$ & $%.4f$ & & & \\\\", 
                data.d[i]*1000, Rc_Mohm
            ))
        end
    end
\end{luacode*}
 

\textit{Темновое сопротивление образца CdS:} $_\ast R_{\text{т}} = \directlua{tex.print(protocol.DarkResistanceOfCdS/1e3)}$ кОм.

\textit{Темновое сопротивление образца CdSe:} $_\ast R_{\text{т}} = \directlua{tex.print(protocol.DarkResistanceOfCdSe/1e6)}$ МОм.

\begin{longtblr}[
    caption = {Спектральная зависимость фотопроводимости (Образец CdS)},
    label = {tab:spectral_cds},
    expand = \directlua,
]{
    width = \linewidth,
    colspec = {X[c] X[c] X[c] X[c] X[c] X[c] X[c] X[c]},
    hlines, vlines,
    rows = {m, font=\small},
    rowhead = 1,
    stretch = 0.75
}
    Дел. & $_{\ast\ast}\lambda$, мкм & $_{\ast\ast}\text{Э}_\lambda$ & $_\ast\!R_\text{с}$, кОм & $\gamma_\text{с}$, мкСм & $\gamma_\text{ф}$, мкСм & $\gamma'_\text{ф}$, у.е. & $\gamma'_\text{ф} / \gamma'_{\text{ф}\max}$ \\
    \directlua{print_spectral_rows("SpectralCdS")}
\end{longtblr}

\begin{longtblr}[
    caption = {Спектральная зависимость фотопроводимости (Образец CdSe)},
    label = {tab:spectral_cdse},
    expand = \directlua,
]{
    width = \linewidth,
    colspec = {X[c] X[c] X[c] X[c] X[c] X[c] X[c] X[c]},
    hlines, vlines,
    rows = {m, font=\small},
    rowhead = 1,
    stretch = 0.75
}
    Дел. & $_{\ast\ast}\lambda$, мкм & $_{\ast\ast}\text{Э}_\lambda$ & $_\ast\!R_\text{с}$, МОм & $\gamma_\text{с}$, мкСм & $\gamma_\text{ф}$, мкСм & $\gamma'_\text{ф}$, у.е. & $\gamma'_\text{ф} / \gamma'_{\text{ф}\max}$ \\
    \directlua{print_spectral_rows("SpectralCdSe")}
\end{longtblr}

\begin{longtblr}[
    caption = {Зависимость фотопроводимости от интенсивности облучения},
    label = {tab:intensity},
    expand = \directlua,
]{
    width = \linewidth,
    colspec = {X[c] X[c] X[c] X[c] X[c]},
    hlines, vlines,
    rows = {m, font=\small},
    rowhead = 1,
    stretch = 0.75
}
    $d$, мм & $_\ast R_\text{с}$, МОм & $\gamma_\text{с}$, мкСм & $\gamma_\text{ф}$, мкСм & $d/d_{\max}$, о.е. \\
    \directlua{print_intensity_rows()}
\end{longtblr}

\ast обязательно для заполнения

\ast\ast заполняется из градуировочной таблицы